\subsection{Exact methods and P1 environment}
\label{subsec:exact-setup}

The Exact family consists of four implementations: Gurobi MIP, CPLEX MIP, CPLEX CP Optimizer, and OR-Tools CP-SAT. All methods receive the same manifest-validated graph instance through the shared project data contract. Each run uses one solver worker. The MIP implementations set zero relative and absolute MIP gaps where those parameters are available; CP Optimizer and CP-SAT use their respective single-worker settings. Symmetry breaking and warm starts are disabled. A seed of zero is supplied to the backends that expose the corresponding setting; the current CPLEX CP adapter does not set a random-seed parameter.

For the official comparison, the current protocol specifies a per-run time limit of 1,800 seconds and a 12-GiB process-level memory cap, with runs executed sequentially on one host. The runner's measured \texttt{wall\_time} is the primary end-to-end runtime. Exact-solver \texttt{t\_encode}, \texttt{t\_solve}, and \texttt{t\_total} are retained as diagnostic fields. The Exact Pilot is a separate operational check over ten validation instances and four Exact backends, with a 600-second limit per run; its observations must not be pooled with, or reported as, the official RQ3 results.

The candidate machine snapshot selected by P1 is Windows 10 build 10.0.26200, an AMD Ryzen 5 7535HS processor, 12 logical CPUs, and 15.21 GiB of physical memory. The selected 12-GiB cap leaves approximately 3.21 GiB outside the benchmark job for the operating system and background processes. The operating system enforces the Windows cap through a Job Object in the shared runner. These values describe P1's candidate environment, not an independently confirmed final official-host declaration; P3 must confirm the official host, resource enforcement, and run accounting before the protocol is frozen.

The recorded package snapshot is OR-Tools 9.15.6755, Gurobi Python API 13.0.3, CPLEX Python API 22.2.0.1, and DOcplex 2.32.264. The Python version must be reconciled before freezing the manuscript: \texttt{docs/w3\_preflight.md} records Python 3.11.16 for the candidate run environment, whereas the older development-environment note \texttt{docs/ENV.md} records Python 3.11.17 in a Conda environment. The exact interpreter used for the official batch must be captured from the selected \texttt{uv} environment and tied to the commit and raw records. The user has not run the 40-run HB Exact Pilot at the time this section was prepared; no pilot outcome is asserted here.

The license preflight is incomplete for two commercial backends. A small Gurobi Academic/WLS probe succeeded, but this does not establish benchmark-scale capacity. The available CPLEX MIP Promotional/Community Edition rejected a probe with 1,001 variables and 1,001 constraints (error 1016); the available CPLEX CP Community Edition rejected a 1,001-binary-variable probe because its search-space limit was exceeded. No CPLEX Academic entitlement has been verified. These limitations are recorded as environment/availability constraints: any affected benchmark attempt must retain its raw error and be classified as \texttt{ERR}, not \texttt{TO}. No license file, key, identifier, or secret is included in the repository.

The 12-GiB value is configured through the shared runner's operating-system memory supervisor; it is not enforced by the Exact model APIs themselves. P1 selected this candidate cap and recorded local/CI memory-limit test evidence, while P3 owns the official-queue operational sign-off. Before the experimental setup is frozen, the group must bind the cap and its termination/status accounting to the actual host and final runner commit. In particular, the official run records must distinguish a confirmed memory-limit termination (\texttt{MO}) from an unexplained process failure (\texttt{ERR}).
